
\subsection{定理 B、C、G 的部分证明}
\label{subsec:3-8-10}

现在汇总迄今得到的结果，并说明它们蕴含定理 B 的 (B.2)、定理 C，
以及定理 G 的 (G.1)、(G.5)（连同定理 B、G 中弗洛尔上同调的构造）。

定理 B 中弗洛尔上同调的构造是命题 3.7.17 与定义 3.8.40 的推论；
(B.2) 随即由命题 3.7.87 (2) 与定义 3.8.38 得出。

定理 C 由定理 3.8.50 与推论 3.8.96 得出。具体地说，除
$m_k<\dim L/2$ 这一限制外，定理 C 的全部陈述都是定理 3.8.50 的
推论。下面利用推论 3.8.96 证明
$[o_k^{\dim L}(L;\mathrm{def})]=0$。置 $n=\dim L$。

考虑定理 3.8.41 的证明。我们将归纳证明，可以选择
$\mathfrak b_{(i)}$ 使
$\Pi^n(\mathfrak b_{(i)})=0$。换言之，$\mathfrak b_{(i)}$ 的
$H^n(M;\mathbb Q)\otimes\Lambda^+_{0,\mathrm{nov}}$ 分量为零。
\translatornote{原书这一释义句只排到“$\mathfrak b_{(i)}$ 的
$H^n(M;\mathbb Q)\otimes\Lambda^+_{0,\mathrm{nov}}$ 分量”便以句号结束，
缺少谓语；紧前的 $\Pi^n(\mathfrak b_{(i)})=0$ 唯一确定所缺内容为
“该分量为零”，译文据此补足。}

假设已经对所有 $i<k$ 选定了这样的 $\mathfrak b_{(i)}$。于是可以
应用推论 3.8.96 (1)。因此，若
\translatornote{原书本段多次把障碍上链
$o_k(L;\mathrm{weak,def})$ 直接写成 $H^n(L;\mathbb Q)$ 中的元素，
并直接令它为零，漏去了表示障碍上同调类的方括号。定理 3.8.50 的
障碍类定义和这里所用的推论 3.8.96 均明确作用于
$[o_k(L;\mathrm{weak,def})]$；译文据此在本段统一补回方括号。}
$[o_k(L;\mathrm{weak,def})]\in H^n(L;\mathbb Q)$，则它自动为零。
这表明，在选择满足 \eqref{eq:3-8-47} 的 $b_k$ 时，只要
$[o_k(L;\mathrm{weak,def})]\in H^n(L;\mathbb Q)$，就可以取
$\mathfrak b_k=0$。注意，$\mathfrak b_k\in H^n(M;\mathbb Q)$ 当且仅当
$[o_k(L;\mathrm{weak,def})]\in H^n(L;\mathbb Q)$。故归纳可以进行。

与此同时，我们还证明了：若
$[o_k(L;\mathrm{weak,def})]\in H^n(L;\mathbb Q)$，则
$[o_k(L;\mathrm{weak,def})]=0$。这正是所要证明的结论。还要注意，
由命题 3.8.54，条件
$[o_k^{\dim L}(L;\mathrm{def})]=0$ 等价于条件
$[o_k^{\dim L}(L;\mathrm{weak,def})]=0$。

定理 G 中弗洛尔上同调的存在性同样是命题 3.7.17 与定义 3.8.58 的
推论。(G.1) 由命题 3.7.75 及其类似结论得出；该类似结论断言，当
$L^{(1)}=L^{(0)}$ 时，算子 $\mathfrak r,\mathfrak r^+$ 分别与
$\mathfrak q,\mathfrak q^+$ 一致。（这一事实的证明与命题 3.7.75 的
证明类似。）最后，(G.5) 是命题 3.7.93 的推论。
